% ============================================================
% Digital Discovery Introduction — ML/AI Emphasis with Application Context
% Flow: Agricultural/human-centric lighting need → TADF opportunity →
%       Computational bottleneck → Existing ML limitations →
%       Our contributions (accelerated discovery framework)
% ============================================================

\section{Introduction}\label{sec:introduction}

\subsection{Application-driven materials discovery: horticultural and human-centric lighting}

Controlled-environment agriculture and human-centric lighting demand spectrally
tunable, energy-efficient light sources, yet commercial inorganic LEDs rely on
rare-earth phosphors (Eu, Tb, Y) with broad emission and supply-chain
risk~\cite{Yokoyama_2018, Bugbee2016}. Thermally activated delayed fluorescence
(TADF) materials are a tunable, heavy-metal-free alternative: organic
donor--acceptor emitters that harvest both singlet and triplet excitons for
near-unity internal quantum efficiency~\cite{Uoyama2012}.
Realizing this potential means navigating a vast donor--acceptor space under
competing requirements---a minimal singlet--triplet gap
($\Delta E_\text{ST} < \qty{0.2}{\electronvolt}$), efficient reverse intersystem
crossing ($k_\text{RISC} > \qty{1e6}{\per\second}$), and spectral overlap with
target photoreceptor bands. This paper addresses only the first,
$\Delta E_\text{ST}$; the rate and device-efficiency requirements demand
spin--orbit-coupling and device-level calculations beyond the cheap structural
data studied here, and we make no quantitative claim about them
(\cref{subsec:limitations}).

\subsection{The computational bottleneck and the role of machine learning}

Designing such emitters means searching a large space of donor--acceptor
combinations for a small $\Delta E_\text{ST}$. The direct route---time-dependent
DFT on excited states---scales steeply with system size~\cite{Curtarolo2013},
so exhaustive screening of more than a few hundred candidates is impractical.
Machine learning offers a cheaper surrogate~\cite{Butler2018, Pyzer-Knapp2022},
and several groups have predicted TADF gaps from molecular
descriptors~\cite{Gomez-Bombarelli2016, Barneschi2024}. Those low reported errors,
however, are obtained against \emph{computed} (TD-DFT or higher) reference gaps;
predicting the \emph{experimental} gap is markedly harder, to the point that one
study abandoned direct regression for classification~\cite{Nigam2024}
(\cref{subsec:comparison}).

Two problems recur in this literature, and we confront both. First, accuracy is
often quoted against a single random train--test split, which rewards
interpolation within familiar scaffolds and overstates how a model behaves on new
chemistry. Second, and more subtle, models are sometimes trained to reproduce a
\emph{computed} gap from a feature set that already contains the constituent
excitation energies; because $\Delta E_\text{ST}\equiv E_{S_1}-E_{T_1}$, this
reduces to subtraction and the reported accuracy reflects arithmetic rather than
learning.

Our earlier work set the stage: Article~1~\cite{Tchapet2024_xTB} benchmarked
GFN2-xTB/sTDA against \num{747} experimental molecules
(MAE = \qty{0.162}{\electronvolt}), and Article~2~\cite{Tchapet2024_DataDriven}
extracted structure--property design rules from that set. Here we ask the
narrower, sharper question those studies left open: starting only from cheap
structural information, how accurately can a model predict the \emph{measured}
$\Delta E_\text{ST}$, and what sets the limit?

\subsection{Our contributions: an honest benchmark of structure-based prediction}

Rather than claim a new high-accuracy predictor, we ask how far cheap,
structure-based machine learning can go in predicting the \emph{experimental}
singlet--triplet gap---and answer carefully, explicit about where it fails. The
primary result is a scaffold-validated benchmark with a calibrated accuracy ceiling:
2D structure alone predicts the measured gap as well as semi-empirical excited-state
features, and we quantify why neither does better. A secondary, cautionary
contribution is methodological---a reminder that excited-state-energy \emph{leakage}
inflates reported accuracy for any property defined as a difference of computed
quantities. We make four contributions.

\paragraph{Structure alone, with no excited-state calculation, does as well as anything.}
On \num{231} molecules across \num{212} Bemis--Murcko scaffolds, random forests built
from the 2D structure alone---Morgan fingerprints (MAE = \qty{0.091}{\electronvolt}) and
RDKit descriptors (\qty{0.100}{\electronvolt})---match descriptors derived from a
semi-empirical excited-state calculation (natural transition orbital, NTO, spatial
features, \qty{0.096}{\electronvolt}). The result holds across scaffold, cluster and random
cross-validation and against a label-permutation null, so it is not a splitting
artefact. The excited-state step buys nothing for the $S_1$--$T_1$ gap; SHAP nonetheless
ties the NTO variant to hole--electron overlap and charge transfer, giving an
interpretable read on what structure encodes.

\paragraph{An explicit, tested accuracy ceiling.}
We quantify why accuracy saturates near \qty{0.1}{\electronvolt}: independent
experimental reports for the same compound disagree by \qty{0.072}{\electronvolt}
(RMS, up to \qty{0.56}{\electronvolt}), and the computed reference is
functional-dependent (B3LYP vs CAM-B3LYP differ by up to
\qty{0.58}{\electronvolt}). The ceiling is not an excuse: no higher-capacity learner
(gradient boosting, MLP, SVR) beats the random forest, and the result places itself
honestly at $R^2 \approx 0.25$, a real but limited signal bounded by data quality
rather than model capacity.

\paragraph{A cautionary note on feature--target leakage.}
Secondary to the benchmark, we flag that regressing $\Delta E_\text{ST}$ on features
containing the $S_1$/$T_1$ energies is circular ($\Delta E_\text{ST} \equiv E_{S_1} -
E_{T_1}$): the model returns arithmetic, invisible to cross-validation. Such
feature--target leakage is a recognised cause of over-optimistic ML results~\cite{Kapoor2023}
and is singled out in best-practice guidance for chemistry,\cite{Artrith2021} yet remains
easy to introduce whenever the target is a derived energy difference.

\paragraph{Modest, quantified triage utility---and no device claims.}
As a pre-screen for TADF viability, the model enriches low-gap candidates
(\numrange{1.2}{1.4}$\times$ over random selection), and we deposit a conformal-bounded
ranked shortlist of the enumerated library---useful for prioritisation but far from
decisive. $\Delta E_\text{ST}$ is necessary but not sufficient for TADF efficiency;
emission-wavelength matching for any specific lighting niche is a separate, downstream
step not addressed here. We make no device-level performance claims.
% ; our earlier
% spin--orbit-coupling estimates of $k_\text{RISC}$ and external quantum
% efficiency proved molecule-independent and were withdrawn.

This study builds on our earlier work---\textbf{Article 1}~\cite{Tchapet2024_xTB}
(xTB/sTDA validation) and \textbf{Article 2}~\cite{Tchapet2024_DataDriven} (design
rules)---by asking, rigorously, what such cheap data can and cannot support.
All data, trained models, and code are openly available (Zenodo DOI:
10.5281/zenodo.17436070), enabling reproduction and extension of the benchmark.
